\section{A Readers and Writers Lock}

We now implement a readers and writers lock using semaphores.  As in
Section~\ref{sec:conditionsReadersWriters}, we aim to prevent starvation of
writers (collectively), and of each reader.

We use the technique of passing the baton.  We use separate semaphores for
signalling (passing the baton) to a reader or to a writer.  A signal on each
of these semaphores indicates that a waiting thread can enter.

The code is in Figures~\ref{fig:SemaphoreReadersWriters1}
and~\ref{fig:SemaphoreReadersWriters2}.  As before, we keep track of the
number of readers and writers in the critical section, using variables
|readers| and |writers|, and maintain the invariant 
\[\mstyle
(\sm{readers} = 0 \lor \sm{writers} = 0) \land \sm{writers} \le 1.
\]
We also keep track of how many readers and writers are waiting, so an
exiting thread can know to whom to pass the baton.  We use a mutex semaphore
to protect the object variables.


%%%%%

\begin{figure}
\begin{scala}
class SemaphoreReadersWriters extends ReadersWritersLock{
  /** Number of readers currently in the critical region. */
  private var readers = 0
  /** Number of writers currently in the critical region. */
  private var writers = 0
  /* Invariant: £$(\sm{readers} = 0 \lor \sm{writers} = 0) \land \sm{writers} \le 1$£. */

  /** Number of readers waiting to enter. */
  private var readersWaiting = 0
  /** Number of writers waiting to enter. */
  private var writersWaiting = 0

  private val mutex = new MutexSemaphore

  /* Invariant: if £writers $= 0$£, £writersWaiting $= 0$£, and £mutex£ is up, then
   * £readersWaiting $= 0$£. */

  /** Semaphore to signal to a waiting reader that it can now enter.  A signal
    * indicates that £writers $= 0$£. */
  private val readerEnterS = new SignallingSemaphore

  /** Condition to signal to a waiting writer that it can now enter.  A signal
    * indicates that £readers $= 0$£ and £writers $= 0$£. */ 
  private val writerEnterS = new SignallingSemaphore

  ... // Continued in Figure £\sl\ref{fig:SemaphoreReadersWriters2}£.
}
\end{scala}
\caption{A readers and writers lock, implemented using semaphores (part~1).}
\label{fig:SemaphoreReadersWriters1}
\end{figure}

%%%%%

As in Section~\ref{sec:conditionsReadersWriters}, a reader trying to enter
waits if there is a waiting writer.  This helps to ensure that writers aren't
starved.  

%%%%%

\begin{figure}
\begin{scala}
  def readerEnter() = {
    mutex.down()
    if(writers > 0 || writersWaiting > 0){ // Thread has to wait.
      readersWaiting += 1
      mutex.up(); readerEnterS.down() // Wait for signal (1).
      assert(writers == 0); readersWaiting -= 1
    }
    readers += 1
    if(readersWaiting > 0) readerEnterS.up() // Signal to next reader at (1).
    else mutex.up()
  }

  def readerExit() = {
    mutex.down()
    assert(writers == 0); readers -= 1 
    if(readers == 0 && writersWaiting > 0) 
      writerEnterS.up() // Signal to writer at (2).
    else{ assert(readersWaiting == 0 || writersWaiting > 0);  mutex.up() } 
  }

  def writerEnter() = {
    mutex.down()
    if(readers > 0 || writers > 0){ // Thread has to wait.
      writersWaiting += 1
      mutex.up(); writerEnterS.down() // Wait for signal (2).
      assert(readers == 0 && writers == 0); writersWaiting -= 1
    }
    writers += 1
    mutex.up()
  }

  def writerExit() = {
    mutex.down()
    assert(writers == 1 && readers == 0); writers -= 1 
    if(readersWaiting > 0) readerEnterS.up() // Signal to reader at (1).
    else if(writersWaiting > 0) writerEnterS.up() // Signal to writer at (2).
    else mutex.up()
  }
}
\end{scala}
\caption{A readers and writers lock, implemented using semaphores (part~2).}
\label{fig:SemaphoreReadersWriters2}
\end{figure}

%%%%%

When a writer exits, it passes the baton to a waiting reader if there is one.
Otherwise, it passes the baton to a waiting writer if there is one.
Otherwise, it raises the mutex, to pass the baton to a thread calling an
operation.  Thus it gives priority to waiting readers, which helps to prevent
readers from being starved. 

When a reader exits, if now $\sm{readers} = 0$, it passes the baton to a
waiting writer if there is one.  Again, this is to help ensure that writers
aren't starved.  It turns out that there cannot be another reader waiting at
this point.  There is an invariant that, when the mutex is up, if
$\sm{writers} = 0$ and $\sm{writersWaiting} = 0$, then $\sm{readersWaiting} =
0$: there is no need for a reader to wait when ${\sm{writers} = 0}$ and
$\sm{writersWaiting} = 0$.  Thus, when the exiting reader acquires the mutex,
if $\sm{writersWaiting} = 0$ there is no waiting reader.  In this case, it
simply raises the mutex to pass the baton to a new operation call.

If a reader has to wait, when it receives a signal (necessarily from a
writer), it passes the baton to another reader if there is one.  Thus, the
baton will be passed from the writer to a reader, and then from one reader to
another, until all the readers have entered.  The final reader simply raises
the mutex. 

\begin{instruction}
Study the details of the implementation.  Persuade yourself that the claimed
invariants hold.
\end{instruction}

%% A leaving reader can signal to a waiting writer.  A leaving writer can signal
%% to either a waiting reader or waiting writer.

%% The previous version could cause a writer to be starved, if readers
%% continually enter and leave the critical section such that there is always at
%% least one reader.  To avoid this, we force a new reader to wait if there is
%% already a waiting writer.

%% On the assumption that no thread remains in the critical section for ever,
%% this ensures that writers collectively are not starved.  Once a writer starts
%% waiting, no more readers will enter the CS; eventually all readers will leave
%% the CS; and a writer will receive a signal.

%% The implementation is potentially unfair to an \emph{individual} writer.  If
%% repeatedly several writers are waiting, a particular writer may never receive
%% the signal.  Likewise, a thread may never receive a signal on |mutex|
%% (although this is less likely). 

%% The implementation of a semaphore does not guarantee to be fair to each
%% individual waiting thread, although in practice it might be.

%%%%%%%%%%%%%%%%%%%%%%%%%%%%%%%%%%%%%%%%%%%%%%%%%%%%%%%

\section{Counting semaphores}

The semaphores we have seen so far have been binary: they have just two
states, \emph{up} and \emph{down}.  By contrast, a \emph{counting semaphore}
has a non-negative integer as its state.  The \SCALA{down} operation waits
until the value is strictly positive, and decrements it.  The \SCALA{up}
operation increments the value, waking a thread if necessary.

Figure~\ref{fig:CountingSemaphore} gives an implementation of a counting
semaphore, using a JVM monitor.  The constructor of the semaphore takes an
optional parameter that represents the initial state, with default value
of~|0|.

%%%%%

\begin{figure}[htbp]
\begin{scala}
class CountingSemaphore(private var permits: Int = 0){
  def down() = synchronized{
    while(permits == 0) wait()
    permits -= 1
  }

  def up() = synchronized{
    permits += 1
    notify()
  }
}
\end{scala}
\caption{The implementation of a counting semaphore.}
\label{fig:CountingSemaphore}
\end{figure}

%%%%%

The value of the semaphore can be thought of as the number of \emph{permits}
available for |down| operations.  Each |down| operation consumes one permit.
Each |up| operation provides one additional permit.

Figure~\ref{fig:CountingSemaphorePartialQueue} gives an implementation of a
partial queue that uses a counting semaphore~|size|, together with a mutex
semaphore that protects the underlying queue.  In particular, the state of
|size| equals the number of items in the queue: it can be thought of as the
number of permits to perform dequeues; initially, it equals the default
value~|0|.  Each |enqueue| operation provides one permit, so performs
|size.up()|.  Each |dequeue| operation consumes one permit, so performs
|size.down()|; this blocks when the queue is empty, as required.

%%%%%

\begin{figure}
\begin{scala}
class CountingSemaphorePartialQueue[T] extends PartialQueue[T]{
  /** The queue itself. */
  private val queue = new scala.collection.mutable.Queue[T]

  /** Semaphore for mutual exclusion on £queue£. */
  private val mutex = new MutexSemaphore()

  /** Semaphore for dequeueing.  The state of the semaphore equals
    * £queue.length£. */
  private val size = new CountingSemaphore(0)

  def enqueue(x: T) = {
    mutex.down()
    queue.enqueue(x)
    size.up()
    mutex.up()
  }

  def dequeue(): T = {
    size.down()
    mutex.down()
    val result = queue.dequeue()
    mutex.up()
    result
  }
}
\end{scala}
\caption{A partial queue implemented using a counting semaphore.}
\label{fig:CountingSemaphorePartialQueue}
\end{figure}


%%%%%

Note that the order in which |dequeue| obtains the semaphores is important: it
performs |size.down()| \emph{before} |mutex.down()|.  If it were to obtain
|mutex| first, and |queue| is empty, then the |size.down()| operation would
block.  But at this point, this thread holds the mutex, so all other threads
would be blocked: the system would be deadlocked!

%% \emph{Exercise:} adapt this example to produce a bounded buffer that cannot
%% hold more than $n$ pieces of data.
